\section{Constructions on Vector Bundles}

\subsection{Review of tensor products}

This section will be rather short, assuming a basic familiarity with tensor products.
The following definition is a brief overview of notation.

\bdefn

    We denote the tensor product of modules by $A\otimes B$.
    We denote the dual of an $R$-module $A$ by $A^*=\hom_R(A,R)$.

\edefn

\bnote

    Primitive elements of the tensor product $A\otimes B$ are written $a\otimes b$; all other elements are sums of these.

    The tensor product is associative (in a natural way), and we denote the tensor of $n$ modules without parentheses.
    In an $n$-fold tensor product, we write the primitive elements without parentheses.

    If $B_i=(b^i_1,\dots,b^i_{n_i})$ is a basis for each $V_1$, then
    $$ B = \set{b^1_{j_1}\otimes b^n_{j_n}}[b^i_{j_i}\in B_i] \subseteq V_1\otimes\cdots\otimes V_n $$
    is a basis for the tensor product.
    In particular the dimension of a tensor product of vector spaces is the product of the dimensions of the vector spaces.

    If $A$ is an $R$-module, then $A\otimes R\cong A$ naturally.
    This natural isomorphism is given by mapping $a\otimes r\mapsto ra$.

    The tensor product is functorial: if $F_i\colon A_i\to B_i$ are maps for $i=1,2$ then define
    $F_1\otimes F_2\colon A_1\otimes A_2\to B_1\otimes B_2$ by
    $$ (F_1\otimes F_2)(a_1\otimes a_2) = (F_1(a_1))\otimes(F_2(a_2)) $$
    and extending by linearity.
    This is well-defined.

    In the case of linear functionals $\phi_i\colon A_i\to\bR$ we have that $\bR\otimes\bR\cong\bR$, and this map is naturally isomorphic to
    $$ (\phi_1\otimes\phi_2)(a_1\otimes a_2) = \phi_1(a_1)\phi_2(a_2) . $$

    The universal property of the tensor product is that
    $$ \hom(V_1\otimes\cdots\otimes V_n,A) \cong L(V_1,\dots,V_n;A) $$
    naturally, where the right side is the space of multilinear maps $V_1\otimes\cdots\otimes V_n\to A$.

\enote

\bnote

    For a field $k$ and a finite-dimensional $k$-linear space $V$, we know that $\dim V^*=\dim V$.
    For a basis $B=(v_1,\dots,v_n)$ we define its \emphasize{dual basis} $B^*=(v_1^*,\dots,v_n^*)$ in $V^*$ as the linear functionals
    defined by
    $$ v_i^*(v_j) = \delta_{ij} , $$
    and extending by linearity.
    These are linearly independent, and thus must be a basis considering dimensions.

    There is a natural map from a space to its double dual, $V\to V^{**}$.
    For $v\in V$ define $\hat v\in V^{**}$ by $\hat v(\phi)=\phi(v)$ for all $\phi\in V^*$.
    This is injective, and considering dimensions is an isomorphism.
    It is natural, and thus for finite-dimensional spaces $V$ and $V^{**}$ are naturally isomorphic.

    For infinite-dimensional spaces, $\dim V^*$ is strictly greater than $\dim V$, so in particular $V$ is naturally a proper subspace of
    $V^{**}$, but they are not isomorphic.

\enote

\bprop

    Let $V_1,\dots,V_n$ be vector spaces, then there is a natural isomorphism
    $$ (V_1\otimes\cdots\otimes V_n)^* \cong V_1^*\otimes\cdots\otimes V_n^* . $$

\eprop

\bproof

    They have the same dimension, and thus we need only find an injective map.
    We define a map $V_1^*\otimes\cdots\otimes V_n^*$ to the dual space on the left by mapping $\phi_1\otimes\cdots\otimes\phi_n$ to the
    literal tensor product of maps:
    $$ (\phi_1\otimes\cdots\otimes\phi_n)(a_1,\dots,a_n) = \phi_1(a_1)\cdots\phi_n(a_n) . $$
    If we choose a dual basis for each $V_i^*$, we see that it maps to a linearly independent set, and thus this is injective.
    \qed

\eproof

In particular, we then have the natural isomorphisms
$$ V_1\otimes\cdots\otimes V_n \cong (V_1\otimes\cdots\otimes V_n)^{**} \cong (V_1^*\otimes\cdots\otimes V_n^*)^* . $$
Thus,
$$ V_1\otimes\cdots\otimes V_n \cong \hom(V_1^*\otimes\cdots\otimes V_n,\bR) \cong L(V_1^*,\dots,V_n^*,\bR) . $$
So we can identify the tensor product $V_1\otimes\cdots\otimes V_n$ with multilinear maps $V_1^*\times\cdots\times V_n^*\to\bR$ naturally.
Similarly, $V_1^*\otimes\cdots\otimes V_n^*$ can naturally be identified with multilinear maps $V_1\times\cdots\times V_n\to\bR$.

\bdefn

    For a space $A$, denote by $T^k(A)$ the $k$-fold tensor product of $A$ with itself:
    $$ T^k(A) = A\otimes\cdots\otimes A . $$
    The empty tensor product is taken to be the ambient ring, so $T^0(A)=R$.

\edefn

\bdefn

    Let $V$ be a finite-dimensional real vector space.
    A \emphasize{covariant $k$-tensor} for $k>0$ is an element of $V^*\otimes\cdots\otimes V^*$.
    Equivalently, it is a multilinear map $V\times\cdots\times V\to\bR$.
    If $\alpha$ is a covariant $k$-tensor for some $k$, it is called a \emphasize{covariant tensor} and its \emphasize{rank} is $k$.
    The space of covariant $k$-tensors is $T^k(V^*)$.

    A \emphasize{contravariant $k$-tensor} for $k>0$ is an element of $V\otimes\cdots\otimes V=T^k(V)$.
    Equivalently, a multilinear map $V^*\times\cdots\times V^*\to\bR$.
    The rank is defined similarly.

    For $k,\ell\neq0$ nonnegative, we define a \emphasize{tensor of type $(k,\ell)$} to be an element of
    $$ T^{(k,\ell)}(V) = V\otimes\cdots\otimes V\otimes V^*\otimes\cdots\otimes V^* = T^k(V) \otimes T^\ell(V^*) . $$

\edefn

Note that
$$ T^{(0,k)}(V) = T^k(V^*),\qquad T^{(\ell,0)}(V) = T^\ell(V) . $$

\bdefn

    We define an action of $S_k$ on $T^k(V)$, viewed as multilinear maps, by
    $$ {}^\sigma\alpha(v_1,\dots,v_n) = \alpha(v_{\sigma(1)},\dots,v_{\sigma(n)}) $$
    for $\sigma\in S_k,\alpha\in T^k(V)$.
    This is indeed a group action, moreover it preserves the linear structure of $T^k(V)$ in that
    $$ {}^\sigma(\alpha+\beta) = {}^\sigma\alpha + {}^\sigma\beta,\qquad {}^\sigma(c\alpha) = c{}^\sigma\alpha . $$

\edefn

Viewing $T^k(V)$ as the $k$-fold tensor product of $V^*$, ${}^\sigma$ acts on primitive elements by
$$ \sigma(v_1\otimes\cdots\otimes v_k) = v_{\sigma^{-1}(1)}\otimes\cdots\otimes v_{\sigma^{-1}(k)} , $$
and extending by linearity.

\bdefn

    A covariant $k$-tensor $\alpha\in T^k(V^*)$ is \emphasize{symmetric} if swapping its inputs does not change its value.
    Equivalently it has a trivial orbit under the above action: ${}^\sigma\alpha=\alpha$ for all $\sigma\in S_k$ (since transpositions
    generate $S_k$).
    Denote the space of symmetric covariant $k$-tensors by $\Sigma^k(V^*)\subseteq T^k(V^*)$.

    Similarly a covariant $k$-tensor $\alpha\in T^k(V^*)$ is \emphasize{antisymmetric} if swapping its inputs swaps the sign of its value.
    Equivalently if ${}^\sigma\alpha=\sgn\sigma\alpha$ for all $\sigma\in S_k$.
    Denote the space of alternating covariant $k$-tensors by $\Lambda^k(V^*)\subseteq T^k(V^*)$.
    Antisymmetric tensors are also called \emphasize{exterior forms} or \emphasize{multi-covectors}.

\edefn

Define the map $\sym\colon T^k(V^*)\to\Sigma^k(V^*)$ by
$$ \sym\alpha = \frac1{k!}\sum_{\sigma\in S_k}{}^\sigma\alpha . $$
This map is called \emphasize{symmetrization}, and it is a projection onto $\Sigma^k(V^*)$.
Indeed,
$$ {}^\tau\sym\alpha = \frac1{k!}\sum_{\sigma\in S_k}{}^{\tau\sigma}\alpha = \sym\alpha $$
so that $\sym$ is well-defined, and $\sym\alpha=\alpha$ iff $\alpha$ is symmetric.
Clearly the equality implies $\alpha$'s symmetry since $\sym\alpha$ is symmetric, and if $\alpha$ is symmetric then
$$ \sym\alpha = \frac1{k!}\sum_{\sigma\in S_k}\alpha = \alpha . $$

\bdefn

    If $\alpha\in\Sigma^k(V^*)$ and $\beta\in\Sigma^\ell(V^*)$ are symmetric, then their tensor $\alpha\otimes\beta$ is not symmetric in general.
    Define their \emphasize{symmetric product} to be the symmetrification of their tensor product:
    $$ \alpha\beta = \sym(\alpha\otimes\beta) . $$
    Explicitly,
    $$ \alpha\beta(v_1,\dots,v_{k+\ell}) =
    \frac1{(k+\ell)!}\sum_{\sigma\in S_{k+\ell}}\alpha(v_{\sigma1},\dots,v_{\sigma k})\beta(v_{\sigma(k+1)},\dots,v_{\sigma\ell}) . $$

\edefn

\bprop

    The symmetric product satisfies the following:
    \benum
        \item it is bilinear,
        \item it is commutative,
        \item it is associative.
    \eenum

\eprop

We will not prove this.

We can also define the \emphasize{alternation} of a tensor by
$$ \alt\alpha = \frac1{k!}\sum_{\sigma\in S_k}\sgn\sigma{}^\sigma\alpha . $$
This is indeed antisymmetric,
$$ {}^\tau\alt\alpha = \frac1{k!}\sum_\sigma\sgn\sigma{}^{\tau\sigma}\alpha = \sgn\tau\alt\alpha . $$
If $\alpha$ is already antisymmetric then
$$ \alt\alpha = \frac1{k!}\sum_\sigma\sgn\sigma{}^\sigma\alpha = \frac1{k!}\sum_\sigma(\sgn\sigma)^2\alpha = \alpha $$
so this is also a projection.

Notice that $\sym\alpha=0$ for antisymmetric tensors of dimension $>1$ and similarly $\alt\alpha=0$ for symmetric tensors of dimension $>1$.

\subsection{Tensor fields}

\bdefn

    Let $M$ be a smooth manifold.
    Define its \emphasize{bundle of covariant $k$-tensors} by
    $$ T^kT^*M = \coprod_{p\in M}T^k(T^*_pM) $$
    where $T_p^*M$ is the dual space of $T_pM$.
    Similarly the \emphasize{bundle of contravariant $k$-tensors} by
    $$ T^kTM = \coprod_{p\in M}T^k(T_pM) , $$
    and in general the \emphasize{bundle of $(k,\ell)$-tensors} by
    $$ T^{(k,\ell)}TM = \coprod_{p\in M}T^{(k,\ell)}(T_pM) . $$

\edefn

To show that $T^{(k,\ell)}TM$ has the unique structure of a vector bundle over $M$, we prove that in general the tensor of two vector bundles is a vector bundle.

\bdefn

    Let $\pi'\colon E'\to M$ and $\pi''\colon E''\to M$ be vector bundles over $M$.
    Define their \emphasize{tensor} to be $E=\coprod_{p\in M}E'_p\otimes E''_p$ with the obvious projection $\pi\colon E\to M$.

\edefn

Let $(\sigma_i)$ and $(\tilde\sigma_i)$ be local frames defined in an open set $U$, and let $\Phi,\tilde\Phi$ be their local trivializations.
Meaning
$$ \Phi(a^i\sigma_i|_p) = (p,\bar a), $$
and by assumption there is a smooth transition function $\tau\colon U\to\gl_n(\bR)$ such that
$$ \tilde\Phi\circ\Phi^{-1}(p,\bar a) = (p,\tau(p)\bar a) . $$
Then if $X=a^i\sigma_i=\tilde a^i\tilde\sigma_i$, we have
$$ \tilde\Phi\circ\Phi^{-1}(p,\bar a) = \tilde\Phi^{-1}(X_p) = (p,\tilde a), $$
so that $\tilde a=\tau(p)\bar a$, i.e.\ $\tilde a^i=\tau(p)^i\bar a$ where $\bul^i$ is the $i$th row.

In particular, $\sigma_i=\tau(p)^j_i\tilde\sigma_j$.

\bprop

    The tensor of two vector bundles has a unique smooth structure making it a vector bundle over $M$ whose fibers are the tensors of the fibers of the original
    vector bundles.

\eprop

\bproof

    Suppose $E',E''$ have rank $k',k''$ respectively.
    Let $(\sigma_i)$ and $(\tau_j)$ be local frames of $E',E''$ respectively over an open subset $U\subseteq M$.
    Then define $\Phi\colon\pi^{-1}U\to U\times\bR^{k'k''}$ by
    $$ \Phi(v^{ij}\sigma_i|_p\otimes\tau_j|_p) = (p,v^{ij}) . $$
    Then for another trivialization $\tilde\Phi$ associated with $(\tilde\sigma_i)$ and $(\tilde\tau_j)$ we have that ($T',T''$ being the transition functions),
    $$ \sigma_i\otimes\tau_j = ((T')^k_i\tilde\sigma_k)\otimes((T'')^\ell_j\tilde\tau_\ell) = (T')^k_i(T'')^\ell_j(\tilde\sigma_k\otimes\tilde\tau_\ell) . $$
    So the transition function is $T^i_j=\sum_{k,\ell}(T')^k_i\cdot(T'')^\ell_j$, which is smooth (as the sum and product of smooth functions) and invertible
    (since $\sigma_i\otimes\tau_j$ and $\tilde\sigma_i\otimes\tilde\tau_j$ form bases).
    \qed

\eproof

We have that
$$ T^{(k,\ell)}TM = \coprod_{p\in M}T^{(k,\ell)}(T_pM) $$
is just the tensor of $TM$ tensored with itself $k$ times and $T^*M$ with itself $\ell$ times, and is thus a vector bundle.

The local frames of $E\otimes E'$ are of the form $(\sigma_i\otimes\tau_j)$ for local frames $\sigma_i$ of $E$ and $\tau_j$ of $E'$.
Thus the local frames of $T^{(k,\ell)}TM$ give representations of the form
$$ A^{i_1\dots i_k}_{j_1\dots j_\ell}\pdv{x^{i_1}}\otimes\pdv{x^{i_k}}\otimes dx^{j_1}\otimes dx^{j_\ell} , $$
and $A^{i_1\dots i_k}_{j_1\dots j_\ell}$ are called the \emphasize{components} of the local frame, we denote them also by $A^I_J$ if the tuples $I,J$ are understood.

We adopt the notation for the sections of the tensor bundles:
$$ \T^k(M) = \Gamma(T^kT^*M),\qquad \T^{(k,\ell)}(M) = \Gamma(T^{(k,\ell)}TM) . $$
We do not give special notation to $\Gamma(T^kTM)$, as it is not as important to us.
Sections of $T^kT^*M$ are called \emphasize{covariant $k$-tensor fields}.
Rough sections of $T^{(k,\ell)}TM$ are maps $A\colon M\to T^{(k,\ell)}TM$ such that for every $p\in M$,
$$ A_p\in T^{(k,\ell)}T_pM = L(T_p^*M,\dots,T_p^*M,T_pM,\dots,T_pM;\ \bR) , $$
so $A_p$ is a multlilinear functional from $k$-times $T_p^*M$ and $\ell$-times $T_pM$.

\bnote

    Note that a tensor of type $(k,\ell)$ is {\it contravariant\/} in its first $k$ arguments and {\it covariant\/} in its last $\ell$.
    This may seem unnatural, and indeed some authors think so!
    Some authors mean a tensor of type $(k,\ell)$ to mean it is {\it co}variant in the first $k$ arguments and {\it contra}variant in its last
    $\ell$; $T^{(k,\ell)}$ may have opposite meanings for different authors!

    The convention adopted here is from Lee, but the course at Hebrew University (2026) utilizes the opposite convention.

\enote

\bprop

    Let $M$ be a manifold, and $A\colon M\to T^{(k,\ell)}TM$ be a rough section.
    Then the following are equivalent:
    \benum
        \item $A$ is smooth,
        \item In every coordinate chart $A$'s components are smooth,
        \item Each point in $M$ is contained in a chart where $A$'s components are smooth,
        \item If $X_1,\dots,X_\ell\in\X(M)$ and $\omega_1,\dots,\omega_k\in\X^*(M)$, the function $A(\bar\omega,\bar X)\colon M\to\bR$ defined by
        $$ A(\bar\omega,\bar X)(p) = A_p(\bar\omega|_p,\bar X|_p) $$
        is smooth.
        \item Whenever $\bar X,\bar\omega$ are smooth vector and covector fields defined in some open subset $U\subseteq M$, $A(\bar\omega,\bar X)$ is smooth on $U$.
    \eenum

\eprop

\bproof
    
    The first three points are clearly equivalent, as they are for general vector bundles.
    We note that if $A=A^I_J\ipdv{x^I}\otimes dx^J$, then
    $$ A\parens{dx^{j_1},\dots,dx^{j_k},\pdv{x^{i_1}},\dots,\pdv{x^{i_\ell}}} = A^I_J , $$
    so that the last point implies the first.
    And if $A$ is smooth, then in terms of local coordinates we have that if $X_i=X_i^j\ipdv{x^j}$ and $\omega^i=\omega^i_jdx^j$ then
    $$ A(\bar\omega,\bar X) = A^I_J\ipdv{x^I}\otimes dx^J(\bar\omega,\bar X) = A^I_J\omega_IX^J, $$
    where $X^J=X_1^{j_1}\cdots X_k^{j_k}$ and similar for $\omega_I$.
    This is smooth, as desired.
    \qed

\eproof

\bprop

    Let $M$ be a manifold, $A\in\T^{(k,\ell)}(M)$, $B\in\T^{(k',\ell')}(M)$, and $f\in C^\infty(M)$.
    Then $A\otimes B$ is a smooth tensor field in $\T^{(k+k',\ell+\ell')}$, and $fA\in\T^{(k,\ell)}$.

\eprop

\bproof

    It is not hard to see that if $A$ has local coordinates $A^I_J$, $B$ has $B^{I'}_{J'}$ then $A\otimes B$ has coordinates
    $$ (A\otimes B)^{II'}_{JJ'} = A^I_JB^{I'}_{J'} , $$
    where we split the tuples $II',JJ'$.
    And similarly
    $$ (fA)^I_J = f\cdot A^I_J. $$
    These are all smooth, as desired.
    \qed

\eproof

Our above proposition showed that a tensor field $A\in\T^{(k,\ell)}(M)$ induces a map
$$ \X^*(M)\times\cdots\times\X^*(M)\times\X(M)\times\cdots\times\X(M)\to C^\infty(M) . $$
This is a multilinear $C^\infty(M)$-module morphism, which can be verified coordinate-wise.

\blemm

    Multilinear $C^\infty(M)$ maps $\X^*(M)^k\times\X(M)^\ell\to C^\infty(M)$ are precisely those induced by $(k,\ell)$-tensor fields.

\elemm

\bproof

    We prove this for $(0,k)$-tensor fields, so $C^\infty(M)$-maps $\X(M)^k\to C^\infty(M)$; the general case is similar, but messier.
    Let $\c A$ be a $C^\infty(M)$-multilinear map.
    We first show that $\c A$ acts locally, so if $X_i$ vanishes in a neighborhood $U$ of $p$ then $0=\c A(\bar X)$.
    Again, let $\psi$ be a bump function supported in $U$ such that $\psi(p)=1$, so that $\psi X_i=0$.
    Then
    $$ 0 = \c A(X_1,\dots,\psi X_i,\dots,X_n)(p) = \psi(p)\c A(\bar X)(p) = \c A(\bar X)(p) $$
    meaning $\c A(\bar X)$ vanishes on $U$.

    Now we show that if $X_i|_p=0$ then $\c A(\bar X)(p)=0$ so that $\c A(\bar X)(p)$ is dependent only on $\bar X$'s values at $p$.
    Let $X_i=X^j_i\ipdv{x^j}$ in a neighborhood of $p$, so the two vector fields agree in a neighborhood and thus we can replace $X$ with this coordinate
    representation and not change $\c A$'s value.
    Then
    $$ \c A(X_1,\dots,X_i,\dots,X_n)(p) = X^j_i(p)\c A(X_1,\dots,\ipdv{x^j},\dots,X_n) = 0 , $$
    as desired.
    Note that we need to extend $X^j_i\ipdv{x^j}$ to a global tensor field, which can be done by the extension lemma.

    Now define $A\colon M\to T^kT^*M$ by
    $$ A_p(v_1,\dots,v_k)(p) = \c A(V_1,\dots,V_k)(p) $$
    where $V_i$ are any vector fields with $V_i(p)=v_i$.
    This is independent of choice of $V_i$, as we showed, and gives the desired tensor field (it is smooth since it induces $\A$ which is smooth).
    \qed

\eproof

Thus we can identify tensor fields with multilinear $C^\infty(M)$ morphisms.

There is a natural isomorphism $\hom(V,W)\cong V^*\otimes W$ given by mapping $\phi\otimes w$ to the linear map
$v\mapsto\phi(v)w$.
To show that this is an isomorphism, note that the dimensions are equal, and that it maps a basis to a linearly independent
set.
In particular, this gives a natural isomorphism
$$ L(V,\dots,V;V) \cong \hom(T^k(V),V) \cong T^k(V^*)\otimes V \cong T^{(1,k)}V . $$
This extends to tensor fields: there is a natural isomorphism of mutlilinear maps $\X(M)^k\to\X(M)$ and multilinear maps
$\X(M)^k\times\X(M)^*\to C^\infty(M)$.
This is because given $\A\colon\X(M)^k\to\X(M)$ we define $\tilde\A\colon\X(M)^k\times\X(M)^*\to C^\infty(M)$ by
$$ \tilde\A(X_1,\dots,X_n,\omega) = \omega(\A(X_1,\dots,X_n)) , $$
recalling that cotangent fields act on vector fields smoothly.
Conversely given $\tilde\A$ we define $\A$ by
$$ \A(X_1,\dots,X_n)\omega = \tilde\A(X_1,\dots,X_n,\omega) $$
recalling that by duality a vector field is just a mutlilinear map $X\colon\X(M)^*\to C^\infty(M)$.

\bprop

    $(1,k)$-tensor fields are naturally isomorphic to $C^\infty(M)$-multilinear maps
    $$ \X(M)\times\cdots\times\X(M)\to\X(M) . $$
    In general $(\ell+1,k)$-tensor fields are naturally isomorphic to $C^\infty(M)$-multilinear maps
    $$ \X^*(M)^\ell\times\X(M)^k \to X(M) . $$

\eprop

This is an important basic result as many tensor fields are given in this manner.

\bexam

    The Lie bracket $[\bul,\bul]$ is a map $\X(M)\times\X(M)\to\X(M)$.
    But it is not multilinear over $C^\infty(M)$, and thus is {\it not\/} a tensor.

\eexam

\bdefn

    A \emphasize{symmetric tensor field} is a covariant tensor field whose image consists of symmetric tensors.

\edefn

\bdefn

    Let $F\colon M\to N$ be smooth, then at any $p\in M$ and any $k$-tensor $\alpha\in T^k(T^*_{Fp}N)$ we define $dF^*_p(\alpha)\in T^k(T^*_pM)$ by
    $$ dF^*_p(\alpha)(v_1,\dots,v_n) = \alpha(dF_p(v_1),\dots,dF_p(v_n)) . $$
    This is the \emphasize{pullback of $F$ at $p$}.
    Then for a covariant $k$-tensor field $A$ on $N$, we define the \emphasize{pullback of $A$ by $F$} to be the rough $k$-tensor field on $M$ given by
    $$ (F^*A)_p = dF^*_p(A_{Fp}) . $$
    This acts on vectors $v_1,\dots,v_k\in T_pM$ by
    $$ (F^*A)_p(v_1,\dots,v_k) = A_{Fp}(dF_p(v_1),\dots,dF_p(v_k)) . $$

\edefn

\bprop

    For $F\colon M\to N$ and $G\colon N\to P$ smooth, $A,B$ covariant $k$-tensor fields on $N$, and $f\in C^\infty(N)$,
    \benum
        \item $F^*(fB)=(f\circ F)F^*B$,
        \item $F^*(df)=d(f\circ F)$,
        \item $F^*(A\otimes B)=F^*A\otimes F^*B$,
        \item $F^*(A+B)=F^*A+F^*B$,
        \item $F^*B$ is continuous, and smooth if $B$ is smooth,
        \item $(G\circ F)^*B=F^*(G^*B)$,
        \item $(\id_N)^*B=B$.
    \eenum

\eprop

\bproof

    Similar to before, by definitions.
    \qed

\eproof

\bcoro

    Let $F\colon M\to N$ be smooth, and $B\in\T^k(N)$.
    If $p\in M$ and $(y^i)$ are smooth coordinates for $N$ in a neighborhood of $F(p)$, then $F^*B$ has the representation in a neighborhood of $p$:
    $$ F^*(B_Idy^I) = (B_I\circ F)d(y^{i_1}\circ F)\otimes\cdots\otimes d(y^{i_k}\circ F) . $$

\ecoro

\bproof

    Immediate from the above proposition.
    \qed

\eproof

\subsection{Lie derivatives of tensor fields}

Suppose $V\in\X(M)$ is a smooth vector field with flow $\theta$ with flow domain $\D$.
Recall that for $t\in\bR$, $\theta_t\colon M_t\to M_{-t}$ is a diffeomorphism with inverse $\theta_{-t}$.
For $p\in M$ and $t$ small enough we have that $(t,p)\in\D$, and thus $\theta_t$ is a diffeomorphism from an open neighborhood of $p$ to an
open neighborhood of $\theta_t(p)$.
So $d(\theta_t)_p^*$ pulls back tensors $T^k(T^*_{\theta(t,p)}M)$ to ones at $p$ by
$$ d(\theta_t)^*_p(A_{\theta_t(p)})(v_1,\dots,v_k) = A_{\theta_t(p)}(d(\theta_t)_pv_1,\dots,d(\theta_t)_pv_n) . $$
This is the value of the pullback tensor field $\theta_t^*A$ at $p$.

\bdefn

    If $V\in\X(M)$ is a smooth vector field, and $A$ a covariant tensor field on $M$, define its \emphasize{Lie derivative} with respect to $V$
    by
    $$ (\L_VA)_p = \evdrv t_{t=0}(\theta_t^*A)_p = \lim_{t\to0}\frac{d(\theta_t)^*_p(A_{\theta(t,p)}) - A_p}t . $$

\edefn

\blemm

    The Lie derivative exists for every $p$ and defines a covariant tensor field on $M$.

\elemm

\bproof

    For $p\in M$ let $(U,(x^i))$ be a smooth chart containing $p$.
    Choose an open interval $0\in J_0$ and an open set $U_0\subseteq U$ such that $\theta$ maps $J_0\times U_0$ into $U$.
    For $(t,x)\in J_0\times U_0$ we can write $\theta$ relative to the coordinates by
    $$ \theta(t,x) = (\theta^1(t,x),\dots,\theta^n(t,x)) . $$
    Then $d(\theta_t)_x\colon T_xM\to T_{\theta(t,x)}$ has matrix representation
    $$ \parens{\pdvof{\theta^i}{x^j}(t,x)} . $$

    If $A$ has coordinate representation
    $$ A = A_Jdx^J , $$
    then
    $$ d(\theta_t)^*_x(A_{\theta_t(p)})\parens{\evpdv{x^I}_x} =
    A_{\theta_t(p)}\parens{\pdvof{\theta^{i_k}}{x^k}(t,x)\evpdv{x^{i_k}}_{\theta(t,x)}} =
    \pdvof{\theta^{i_1}}{x^1}(t,x)\cdots\pdvof{\theta^{i_k}}{x^k}(t,x)\cdot A_{i_1,\dots,i_k} . $$
    This is a smooth coordinate representation and thus has a smooth derivative.
    \qed

\eproof

\bprop

    Let $V\in\X(M)$ be smooth, $f\in C^\infty(M)$ (which can be regarded as $T^0TM$), and $A,B$ are smooth covariant tensor fields on $M$
    (of potentially different type).
    \benum
        \item $\L_Vf=Vf$,
        \item $\L_V(fA)=(\L_Vf)A+f(\L_VA)$,
        \item $\L_V(A\otimes B)=(\L_VA)\otimes B+A\otimes(\L_VB)$,
        \item if $X_1,\dots,X_k$ are smooth vector fields and $A$ has type $k$, then
        $$ \L_V(A(X_1,\dots,X_k)) = (\L_VA)(X_1,\dots,X_k) + A(\L_VX_1,\dots,X_k) + \cdots + A(X_1,\dots,\L_VX_k) . $$
        Since $\L_VX=[V,X]$ for vector fields $X$, we have (by the first point the left-hand side above is just $V$ on $A$'s value),
        $$ (\L_VA)(X_1,\dots,X_n) = V(A(X_1,\dots,X_n)) - A([V,X_1],\dots,X_k) - \cdots - A(X_1,\dots,[V,X_k]) . $$
    \eenum

\eprop

\bproof

    Let $\theta$ be $V$'s flow, then
    $$ \theta_t^*(f(p)) = f(\theta_t(p)) = f\circ\theta^{(p)}(t) . $$
    So the Lie derivative $\L_Vf$ is the derivative of this:
    $$ (\L_Vf)(p) = \evdrv t_{t=0}f\circ\theta^{(p)} = df_p({\theta^{(p)}}'(0)) = df_p(V_p) = Vf(p) . $$

    For the other points, if $p$ is a regular point then we choose a coordinate domain where $V$ has representation $\ipdv{u^1}$.
    Then $V$'s flow is just $\theta_t(u)=(u^1+t,u^2,\dots,u^n)$ and then in these coordinates if $A=A_Jdx^J$ we have that $\L_VA$'s coordinate
    representation is just $A_J$'s partial derivatives with respect to $u^1$.
    Then all the other points follow by the product rule.

    Then for points in $V$s support it follows by continuity, and for points outside of the support everything becomes $0$.
    \qed

\eproof

For $f\in C^\infty(M)$ we have the covariant $1$-tensor $df$ given by $df(X)=Xf$.
Notice then that
$$ (\L_Vdf)(X) = V(df(X)) - df([V,X]) = V(Xf) - [V,X]f = VXf - VXf + XVf = XVf , $$
and this is just $df(XV)$ which is equal to $d(Vf)(X)$ and by above this is equal to $d(\L_Vf)(X)$.
Thus we have
$$ \L_Vdf = d(\L_Vf) , $$
for all $f\in C^\infty(M)$.

\bdefn

    If $A$ is a covariant tensor field on $M$, and $\theta$ is a flow on $M$, then $A$ is \emphasize{invariant under $\theta$} if for all
    $(t,p)$ in $\theta$'s domain
    $$ d(\theta_t)^*_p(A_{\theta_t(p)}) = A_p . $$
    Equivalently we may write
    $$ (\theta_t)^*A = A . $$

\edefn

\bthrm

    $A$ is invariant under $V$'s flow iff $\L_VA=0$.

\ethrm

\bproof

    Let $V$'s flow be $\theta$.
    Suppose that $A$ is invariant under $\theta$, then by definition
    $$ \L_V(A)_p = \lim_{t\to0}\frac{d(\theta_t)^*_p(A_{\theta_t(p)}) - A_p}t = \lim_{t\to0}\frac{A_p-A_p}t = 0 $$
    as desired.

    Conversely if $\L_V(A)=0$ then the derivative of $(\theta_t^*A)_p$ at time $0$ is zero.
    Now let $t_0\in\bR$ and notice
    $$ \evdrv t_{t=t_0}(\theta_t^*A) = \evdrv t_{t=0}(\theta_{t+t_0}^*A) = \evdrv t_{t=0}(\theta_{t_0})^*(\theta_t)^*A
    = (\theta_{t_0})^*\evdrv t_{t=0}(\theta_t)^*A = (\theta_{t_0})^*\L_V(A) . $$
    Thus we get that
    $$ \evdrv t_{t=t_0}(\theta_t^*A) = 0 $$
    meaning $\theta_t^*A$ is constant in $t$, so $\theta_t^*A=\theta_0^*A=A$, because $\theta_0=\id$, and thus $A$ is
    invariant under $\theta$ as desired.
    \qed

\eproof

\subsection{The Serre-Swan theorem}

In this section we investigate the category of vector bundles over a smooth manifold $M$.
Recall the strong connection between the $C^\infty(M)$-module of sections of a vector bundle and the vector bundle itself.
In particular we recall the following result we proved earlier.

\bthrm

    Let $E\to M,F\to M$ be vector bundles over $M$.
    There is a natural correspondence between vector bundle morphisms $E\to F$ over $M$ and $C^\infty(M)$-module morphisms $\Gamma(E)\to\Gamma(F)$.

\ethrm

\bproof

    We have already proven this, but we restate the construction of the bijections.
    Given a vector bundle morphism $\phi\colon E\to F$, we define $\tilde\phi\colon\Gamma(E)\to\Gamma(F)$ by $s\mapsto\phi\circ s$.
    Conversely given an module morphism $\psi\colon\Gamma(E)\to\Gamma(F)$ define $\hat\psi\colon E\to F$ by $\hat\psi(v_p)=\psi(\hat v)_p$ where $\hat v$ is any
    section whose value at $p$ is $v$.

    Standard arguments are used to show that $\hat\psi$ is well-defined: first one shows $\psi$ acts locally, so sections agreeing on an open neighborhood
    have the same value under $\psi$ in that neighborhood, then one shows that sections agreeing at a point agree at that point under $\psi$.
    Then $\hat\psi$ is shown to be smooth by considering its representation under local trivializations.

    It is clear that $\tilde\bul$ and $\hat\bul$ are inverses.
    \qed

\eproof

Let us denote by $\VBun_M$ the category of vector bundles over $M$.
Then we have a functor
$$ \Gamma\colon \VBun_M \to \Mod_{C^\infty(M)},\qquad E\mapsto\Gamma(E) , $$
whose actions on vector bundle morphisms is described in the above theorem: $\Gamma(\phi)(s)=\phi\circ s$.
The theorem says that $\Gamma$ is fully faithful, and so we lose no information identifying $\phi$ and $\Gamma(\phi)$.

Because $\Gamma$ is fully faithful, it is an equivalence onto its essential image.
As we will show later, its essential image is the collection of {\it projective\/} $C^\infty(M)$-modules.

\bprop

    A vector bundle $E$ is trivial iff $E$ has a global frame iff $\Gamma(E)$ is free.

    Conversely any free $C^\infty(M)$-module of finite rank is isomorphic to the sections of some trivial vector bundle.

\eprop

\bproof

    The first equivalence was already shown, as the local trivialization related to a global frame is an isomorphism to a trivial bundle.
    And a basis for $\Gamma(E)$ is precisely a global frame for $E$.

    Suppose $F$ is a free $C^\infty(M)$-module, generated by $f_1,\dots,f_n$.
    Then $\Gamma(M\times\bR^n)$ is free of rank $n$ by the previous paragraph, and thus isomorphic to $F$.
    \qed

\eproof

For a bundle morphism $\phi\colon E\to F$ and $p\in M$ let us denote by $\phi_p=\phi|_{E_p}\colon E_p\to F_p$, which is linear by definition.

Now recall the following result we proved as well.

\bprop

    A morphism of vector bundles over $M$ is an isomorphism iff it is bijective.

\eprop

\bproof

    Suppose $\pi\colon E\to M$ and $\pi'\colon F\to M$ are vector bundles and $\phi\colon E\to F$ is a bijective vector bundle morphism.
    Since $\phi_p\colon E_p\to F_p$ is then a linear isomorphism, $E$ and $F$ have the same rank and thus dimension.
    Since $\phi$ is by assumption bijective, it is sufficient to show that it is a local diffeomorphism, and for that it is sufficient to show that it is
    an immersion.

    Now, we have
    $$ \pi = \pi'\circ\phi, $$
    so for all $v\in E$, say $v\in E_p$, we have
    $$ d\pi_v = d\pi'_{\phi(v)}\circ d\phi_v . $$
    Let $x\in T_vE$ such that $d\phi_v(x)=0$ and thus $d\pi_v(x)=0$.

    Recall that for a submersion $\pi$, $v\in\pi^{-1}p$, we have that $\ker d\pi_v=T_v\pi^{-1}p$ (considering dimensions they are equal, and $\pi$ is constant
    on $\pi^{-1}p$ so its differential is zero).
    Thus $x\in T_vE_p$, but $\phi_p$ is an isomorphism $E_p\to F_p$ and thus $d\phi$ is an isomorphism on $T_vE_p$, so $x=0$.
    Thus $d\phi_v$ is an injection, so $\phi$ is a diffeomorphism.
    \qed

\eproof

Another way we could phrase this is that vector bundles are isomorphisms iff their actions on fibers are.

Recall the following result which we proved.

\blemm

    For a vector bundle morphism $\phi\colon E\to F$, the following are equivalent:
    \benum
        \item $\phi_p$ has constant rank for all $p\in M$,
        \item $\ker\phi$ is a subbundle of $E$,
        \item $\im\phi$ is a subbundle of $F$.
    \eenum

\elemm

For a point $p\in M$ let us consider the map $\ev_p\colon C^\infty(M)\to\bR$ which evaluates a smooth function at $p$.
This map is clearly surjective (considering constant functions), and thus its kernel is a maximal ideal.
We denote its kernel by
$$ \mu_p = \ker\ev_p = \set{f\in C^\infty(M)}[f(p)=0] . $$

\blemm

    For $s\in\Gamma(E)$, if $s(p)=0$ then there are $f_i\in\mu_p$ and $s_i\in\Gamma(E)$ such that $s=\sum_if_is_i$.
    In other words, if $s(p)=0$ then $s\in\mu_p\Gamma(E)$ (the product of a module and an ideal).

    Therefore, for all $p\in M$ the following sequence splits
    $$ 0 \to \mu_p\Gamma(E) \to \Gamma(E) \to E_p \to 0 , $$
    where the first map is inclusion, and the second is evaluating a section at $p$.
    Thus we have
    $$ E_p \cong \frac{\Gamma(E)}{\mu_p\Gamma(E)} . $$

\elemm

\bproof

    Let $U$ be a neighborhood of $p$ over which we have a local frame $e_1,\dots,e_n$.
    So there are smooth functions $f_i\in C^\infty(U)$ such that
    $$ s = \sum_i f_ie_i $$
    over $U$, and in particular we have $f_i(p)=0$.
    Now let $\psi$ be a bump function centered at $p$ supported in $U$, so $\psi(p)=1$ and $\supp\psi\subseteq U$.
    Then we have
    $$ \psi^2 s = \sum_i\psi^2 f_ie_i = \sum_i(\psi f_i)(\psi e_i) . $$
    Note that $\psi f_i$ and $\psi e_i$ can be extended smoothly over all of $M$ by setting them to be zero outside of $U$.
    Then this equation holds over all of $M$, and we have
    $$ s = (1-\psi^2)s + \sum_i(\psi f_i)\psi e_i . $$
    And clearly $1-\psi^2,\psi f_i\in\mu_p$ as desired.

    The map $\Gamma(E)\to E_p$ is surjective as for every $v\in E_p$ there is a global section whose value at $p$ is $v$ (which we proved earlier).
    The kernel of the map $\Gamma(E)\to E_p$ is the set of all sections $s\in\Gamma(E)$ whose value at $p$ is zero.
    We showed that this is equal to $\mu_p\Gamma(E)$, showing exactness.
    \qed

\eproof

\blemm

    Let $E,F$ be vector bundles over $M$.
    Then $\Gamma(E\oplus F)\cong\Gamma(E)\oplus\Gamma(F)$.

\elemm

\bproof

    If $s\in\Gamma(E)$ and $t\in\Gamma(F)$ then $s\oplus t\in\Gamma(E\oplus F)$ where $(s\oplus t)|_p=s|_p\oplus t|_p$.
    It is easily seen via local frames that $s\oplus t$ is smooth, and that this map is an injective morphism.
    Similarly for $s\in\Gamma(E\oplus F)$ we map it to $s_0\oplus s_1\in\Gamma(E)\oplus\Gamma(F)$ where $s_0|_p$ is the projection of $s|_p$ onto $E_p$
    and similar for $s_1$.
    Again this is smooth, and these define inverse morphisms.
    \qed

\eproof

Now we state the following result, which will not be proven.

\blemm

    For any vector bundle $E$, $\Gamma(E)$ is a finitely generated $C^\infty(M)$-module.

\elemm

\blemm

    For any vector bundle $E$, $E$ is a direct summand of a trivial vector bundle.

\elemm

\bproof

    Since $\Gamma(E)$ is fg, let $s_1,\dots,s_n\in\Gamma(E)$ generate it.
    Let $Q$ be the free rank-$n$ $C^\infty(M)$-module, and define the surjection $\phi\colon Q\to\Gamma(E)$ by mapping $Q$'s basis to the $s_i$.
    Since $Q$ is free, it is isomorphic to $\Gamma(F)$ where $F$ is a free vector bundle, $F=M\times\bR^n$.
    Thus $\phi\colon\Gamma(F)\to\Gamma(E)$ is a morphism, and thus can be identified with a bundle morphism $\phi\colon F\to E$.

    Since $\Gamma(\phi)$ is surjective, $\phi_p$ is surjective for all $p$.
    This is because we have the following commutative diagram

    \commsq{\Gamma(E)&E_p\cr \Gamma(F)&F_p}
        {}{}{\Gamma(\phi)}{\phi_p}

    and all the arrows other than $\phi_p$ are surjective, implying $\phi_p$ is surjective.
    Thus $\phi$ has constant rank, and $\ker\phi\subseteq F$ is a subbundle.

    Now we claim that $\ker\phi$ is a direct summand of $F$.
    Let $F=M\times\bR^n$, then let $\ker\phi^\perp$ be the subbundle of $F$ whose fiber at $p\in M$ is the orthogonal complement of $\ker\phi_p$ in $\bR^n$
    under the Euclidean product.

    Now the Euclidean product of sections of $F$, which defines a function $\iprod{s,t}_p=\iprod{s(p),t(p)}$, is smooth.
    This is because the Euclidean product of the standard global frames $e_i$ are constant, and the Euclidean product is tensorial, so smoothness follows from linearity.
    If we have a local frame for $\ker\phi$ then extending it to a local frame for $F$ and then applying Gram-Schmidt at each point gives us a local frame for $\ker\phi^\perp$.
    This is because Gram-Schmidt is obtained via the Euclidean product which is smooth.
    Thus $\ker\phi^\perp$ is a subbundle, and clearly $F=\ker\phi\oplus\ker\phi^\perp$.

    Now we claim that the surjection $\phi\colon F\to E$ induces an isomorphism $\phi\colon\ker\phi^\perp\to E$.
    This is because at each $p\in M$, $\phi_p$'s kernel is $\ker\phi_p=(\ker\phi)_p$ and thus forms an isomorphism $\ker\phi^\perp_p\to E_p$.
    So $\phi\colon\ker\phi^\perp\to E$ is an isomorphism on the fibers, and thus an isomorphism of vector bundles, as desired.
    \qed

\eproof

\bcoro

    For a vector bundle $E$, $\Gamma(E)$ is a projective $C^\infty(M)$-module.

\ecoro

\bproof

    $E$ is a direct summand of a trivial vector bundle, so $E\oplus F$ is trivial for some vector bundle $F$.
    Then we have that $\Gamma(E\oplus F)\cong\Gamma(E)\oplus\Gamma(F)$ is free, so $\Gamma(E)$ is projective.
    \qed

\eproof

\bthrm[title=Serre-Swan]

    Let $M$ be a connected manifold.
    If $P$ is a finitely-generated projective $C^\infty(M)$-module, it is isomorphic to the space of sections for some vector bundle $E$.

    Thus $\Gamma$ forms an equivalence of categories from $\VBun_M$ to the full subcategory of finitely-generated projective $C^\infty(M)$-modules.

\ethrm

\bproof

    Since $P$ is projective, it is the direct summand of some finite-rank (since it is finitely-generated) free module.
    So we have $\Gamma(F)=P'\oplus Q$ where $P',Q$ are modules and $P'\cong P$.
    We identify $P$ with $P'$, and thus we can view $P'$ as a set of sections of $F$.

    For $x\in M$ let us define
    $$ P_x = \set{p(x)}[p\in P] \subseteq F_x . $$
    This is a subspace of $F_x$, and similarly we define $Q_x$.
    Now we claim $F_x=P_x\oplus Q_x$.

    For $v\in F_x$ there is a $s\in\Gamma(F)$ such that $s(x)=v$.
    Writing $s=p+q$ for $p\in P,q\in Q$, we have $v=p(x)+q(x)$, showing $P_x+Q_x=F_x$.
    Now if $p(x)=q(x)$ for $p\in P$ and $q\in Q$, then $(p-q)(x)=0$ so $p-q\in\mu_x\Gamma(F)$ and thus
    $$ p-q = \sum_if_is_i = \sum_if_ip_i + \sum_if_iq_i, $$
    where $f_i\in\mu_x$ and $s_i\in\Gamma(F)$ equal to $s_i=p_i+q_i$ for $p_i\in P,q_i\in Q$.
    But then $\sum_if_ip_i\in P$ and so we have $p=\sum_if_ip_i$ and $q=-\sum_if_iq_i$.
    Thus $p(x)=\sum_if_i(x)p_i(x)=0$ and similarly $q(x)=0$, so $v=0$.
    So we have shown $F_x=P_x\oplus Q_x$ as desired.

    Now we want to show that the union of $P_x$ forms a subbundle of $F$, and that $P$ gives its space of sections.
    Notice that $\dim P_x$ is independent of $x$.
    This is because if we take a basis for $P_x$ of $p_1(x),\dots,p_n(x)$, then by continuity $p_1(y),\dots,p_n(y)$ remains linearly independent
    in a neighborhood of $x$.
    Thus $\dim P_y\geq\dim P_x$ in a neighborhood of $x$.
    Similarly we see $\dim Q_y\geq\dim Q_y$, and since $Q_x\oplus P_x=F_x$ whose dimension is constant, we see $\dim P_y$'s is constant in this neighborhood.
    Because $M$ is connected, $\dim P_x$ must be constant over all of $M$.

    If we choose any basis of $P_x$ by $p_1(x),\dots,p_n(x)$, then by continuity $p_1,\dots,p_n$ are linearly independent in a neighborhood of $x$.
    Since we just showed that $\dim P_x$ is constant, this means $p_1,\dots,p_n$ forms a local frame of $P$ in a neighborhood of $x$.
    This means that $\bar P=\coprod_xP_x$ is indeed a subbundle of $F$.

    Now we want to show that $P=\Gamma(\bar P)$.
    Clearly $P\subseteq\Gamma(\bar P)$ since $P$ consists of sections of $F$ falling into $\bar P$.
    And conversely for $s\in\Gamma(\bar P)\subseteq\Gamma(F)$ we have $s=p+q$ for $p\in P$ and $q\in Q$.
    Now $s(x)=p(x)+q(x)$ and since $s(x)\in P_x$, we must have $q(x)=0$.
    So $s=p\in P$ as desired.

    Since $\Gamma$ is fully faithful, it is an equivalence onto its essential image.
    The previous corollary showed that $\Gamma$'s essential image is contained in the subcategory of projective modules, and this one showed that it is in fact
    equal to it.
    \qed

\eproof

\subsection{Functorial constructions on vector bundles}

Thus far we have introduced some methods of constructing vector bundles from old ones.
For example the Whitney sum, the dual bundle, and the tensor product of bundles.
These are all liftings of constructions on vector spaces (direct sums, dual spaces, and tensor products respectively).
In fact, we can generalize this as follows.

\bdefn

    Let $\Vec=\Vec_\bR^{\rm fd}$ denote the category of finite-dimensional vector spaces over $\bR$.
    Then a single-variable covariant \emphasize{$C^\infty$ functor} is a functor $\F\colon\Vec\to\Vec$ such that the maps $\hom_\bR(V,W)\to\hom_\bR(\F(V),\F(W))$ are smooth
    (both sides are vector spaces).
    Contravariant $C^\infty$ functors are defined similarly, and we can extend this to multivariate functors of mixed variance.

\edefn

A multivariate $C^\infty$ functor $\F(x_1,\dots,x_n)$ is smooth iff $\F(V_1,\dots,x_i,\dots,V_n)$ is a single-variable $C^\infty$ functor.

Examples of $C^\infty$ functors are direct sums, taking the dual, and tensor products, as their actions on the hom-sets are linear and thus smooth.
Then the construction on vector bundles done thus far is a special case of the following theorem.

Recall that a local trivialization is a diffeomorphism $\Phi\colon\pi^{-1}U\to U\times\bR^n$ for an open subset $U\subseteq M$
such that projecting onto $U$ gives $\pi$.
So we can write $\Phi$ by $(\pi,\Phi_0)$, and $\Phi_0$ restricts to linear isomorphisms $E_p\to\bR^n$ on each $E_p$.
In the following theorem we will identify a local trivialization with its second component ($\Phi_0$), as no information is lost.

\bthrm

    Let $\F$ be a $C^\infty$ functor.
    For vector bundles $E_1,\dots,E_n\to M$ we define the vector bundle $\F(E_1,\dots,E_n)$ by
    $$ \F(E_1,\dots,E_n) = \coprod_{p\in M}\F(E_1|_p,\dots,E_n|_p) , $$
    with the obvious projection, and this is indeed a vector bundle over $M$.
    Furthermore this construction is functorial.

\ethrm

\bproof

    Suppose $\F$ is covariant in the first $n$ variables and contravariant in the last $m$.
    Let us denote by $E_1,\dots,E_n$ its first $n$ inputs and $F_1,\dots,F_m$ its last $m$.
    If we have local trivializations $\Phi_i$ of the $E_i$ and $\Psi_i$ of the $F_i$ all defined on the same open subset (by potentially shrinking it), we define
    a local trivialization $\Theta$ of $\F(\bar E,\bar F)$ as
    $$ \Theta_p = \F(\Phi_1|_p,\dots,\Phi_n|_p,\Psi_1|_p^{-1},\dots,\Psi_m|_p^{-1}) . $$
    That is, this is the definition of $\Theta$ on the fiber of $p$.
    Since $\F$ is contravariant in its last $m$ arguments, this defines a map from $\F(\bar E_p,\bar F_p)$ to $V=\F(\bR^{N_1},\dots,\bR^{N_n},\bR^{M_1},\dots,\bR^{M_m})$.
    Furthermore it is an isomorphism as functors preserve isomorphisms.

    Now we must just show that different $\Theta$s are compatible.
    So suppose we have other local trivializations $\Phi_i'$ and $\Psi'_i$ giving us $\Theta'$.
    Then
    $$ \Theta'\circ\Theta^{-1}|_p = \Theta'|_p\circ\Theta|_p^{-1} = \F(\bar\Phi'|_p,(\bar\Psi')|_p^{-1})\circ\F(\bar\Phi|_p,\bar\Psi|_p^{-1})^{-1}
    = \F(\bar{\Phi'\circ\Phi^{-1}}|_p,\bar{\Psi\circ(\Psi')^{-1}}|_p) . $$
    In the final equality we swap the order of $\Psi$ and $\Psi'$ due to contravariance.
    So if we let $\tau_i$ be the transition map from $\Phi_i$ to $\Phi'_i$, and $\sigma_i$ the transition from $\Psi'_i$ to $\Psi_i$, we see that the transition map
    $\theta$ from $\Theta$ to $\Theta'$ is given by
    $$ \theta(p) = \F(\bar\tau(p),\bar\sigma(p)) . $$
    Since $\F$ is smooth, we see that $\theta$ is smooth as desired.

    Now suppose we have vector bundle morphisms $\phi_i\colon E_i\to E_i'$ and $\psi_i\colon F'_i\to F_i$.
    Then define $\chi\colon\F(\bar E,\bar F)\to\F(\bar E',\bar F')$ defined on fibers by
    $$ \chi_p = \F(\bar\phi_p,\bar\psi_p) . $$
    If we let $\Theta$ be a trivialization of $\F(\bar E,\bar F)$ and $\Theta'$ a trivialization of $\F(\bar E',\bar F')$ we see that
    $$ \Theta'\circ\chi\Theta^{-1}|_p = \F(\bar\Phi'\circ\bar\phi\circ\bar\Phi^{-1}|_p,\ \bar\Psi\circ\bar\psi\circ(\bar\Psi')^{-1}|_p) . $$
    Since $\phi_i,\psi_i$ are smooth, the arguments of $\F$ on the right are smooth in $p$ and thus the left is smooth in $p$ as well.
    Thus $\chi=\F(\bar\phi,\bar\psi)$ is smooth.
    This is clearly functorial.
    \qed

\eproof

So for every $C^\infty$ functor we associate a functor of vector bundles.
This association itself is functorial: natural transformations of $C^\infty$ functors map to natural transformations of functors of vector bundles functorally.
In particular this means that isomorphic $C^\infty$ functors give isomorphic vector bundle functors.
This allows us to prove easily that the double dual of any vector bundle is naturally isomorphic to the original vector bundle for example.

Or consider the hom functor, $\hom_\bR\colon\Vec^\op\times\Vec\to\Vec$.
This is a linear functor (as in its actions on morphisms is linear) and thus smooth.
Thus it defines a hom functor on vector bundles: for two vector bundles $E,F$ we have a vector bundle $\hom(E,F)$ whose fibers are $\hom_\bR(E_p,F_p)$.
Since we have a natural isomorphism
$$ \hom_\bR(V,U) \cong V^*\otimes U, $$
this lifts to a natural isomorphism
$$ \hom(E,F) \cong E^*\otimes F . $$

Now it is natural to ask what we can say about the space of sections of a functorial construction such as this.
For example, we already know $\Gamma(E\oplus F)\cong\Gamma(E)\oplus\Gamma(F)$.
It is not hard to see that $\Gamma(E^*)\cong\Gamma(E)^*$, where the latter dual is the dual over $C^\infty(M)$: $\Gamma(E)^*=\hom(\Gamma(E),C^\infty(M))$.

To do this we will first prove a few facts about projective modules.

\blemm

    Let $A$ be a commutative ring, and $P,Q$ be finitely generated projective $A$-modules.
    Then
    \benum
        \item $P\oplus Q$ is projective and $(P\oplus Q)^*\cong P^*\oplus Q^*$.
        \item $P$ is \emphasize{reflexive}: $P^{**}\cong P$.
        \item $P\otimes Q$ is projective.
        \item $P^*\otimes Q\cong\hom_A(P,Q)$ and $(P\otimes Q)^*\cong P^*\otimes Q^*$.
        \item $\hom_A(P,Q)$ is projective, in particular $P^*$ is.
    \eenum

\elemm

\bproof

    \benum
        \item $P\oplus Q$ is clearly projective.
        We have
        $$ (P\oplus Q)^* = \hom_A(P\oplus Q,A) \cong \hom_A(P,A) \oplus \hom_A(Q,A) = P^* \cong Q^* , $$
        as desired and this holds for arbitrary modules.

        \item There is a natural map $P\to P^{**}$ which maps $x\in P$ to $\bar x\in P^{**}$ which maps $\phi\in P^*$ to $\phi(x)$.
        For fg free $F$ this map is an isomorphism, as $F^*$ has the same rank as $F$ and it is easy to see that this map maps a basis of $F$ to a linearly independent set of $F^*$.
        So if $P\oplus Q$ is free, then we have an isomorphism $P\oplus Q\to P^{**}\oplus Q^{**}$.
        It is easy to see that this isomorphism is the direct sum of the maps $P\to P^{**}$ and $Q\to Q^{**}$, so they must be isomorphisms.

        \item The tensor of two free modules is free, so if $P\oplus P',Q\oplus Q'$ are free then so is
        $$ (P\oplus P')\otimes(Q\oplus Q') \cong P\otimes Q \oplus \dots . $$
        So $P\otimes Q$ is projective.

        \item There is a natural map $P^*\otimes Q\to\hom_A(P,Q)$ and this is an isomorphism when $P$ is fg projective (independent of $Q$).
        The map is given by mapping $\phi\otimes y$ to $x\mapsto \phi(x)y$.
        For finite free $P$ this is an isomorphism: take $x_1,\dots,x_n$ a basis then every $f\colon P\to Q$ is determined by $f(x_1),\dots,f(x_n)$ and we have a dual basis
        $\phi_1,\dots,\phi_n$ in $P^*$ given by $\phi_i(x_j)=\delta_{ij}$.
        Then $\sum_i\phi_i\otimes f(x_i)$ maps to $f$.

        Now if $P\oplus P'$ is free then we have an isomorphism
        $$ (P^*\otimes Q) \oplus ((P')^*\otimes Q) \to \hom_A(P,Q) \oplus \hom_A(P',Q) $$
        and as before it is easy to see that this is the direct sum of the maps $P^*\otimes Q\to\hom_A(P,Q)$ so they must be isomorphisms.

        \item Suppose $P\oplus P'$ and $Q\oplus Q'$ are free of finite rank.
        Then
        $$ \hom_A(P\oplus P',Q\oplus Q') \cong \hom_A(P,Q) \oplus \hom_A(P,Q') \oplus \hom_A(P',Q) \oplus \hom_A(P',Q') . $$
        And the left-hand side is free, since finite biproducts commute with $\hom$, so $\hom_A(P,Q)$ is a direct summand of a free module and thus projective.

        Since $P^*=\hom_A(P,A)$ and $A$ is free and thus projective, clearly $P^*$ is projective.
        \qed
    \eenum

\eproof

\bprop

    For a vector bundle $E$, $\Gamma(E^*)\cong\Gamma(E)^*$.

\eprop

\bproof

    Define $\phi\colon\Gamma(E^*)\to\Gamma(E)^*$ by $\phi(\bar s)(t)=\bar s(t)$, where $\bar s(t)\in C^\infty(M)$ is given by $\bar s(t)(p)=\bar s_p(t_p)$.
    It is easily seen that $\bar s(t)$ is smooth by considering trivializations.

    For trivial $P$, global sections of $P$ yield global sections of $P^*$, so $P^*$ is trivial as well and so we have an isomorphism $\Gamma(P^*)\cong\Gamma(P)\cong\Gamma(P)^*$
    of free modules, obtained by mapping the dual global sections of $\Gamma(P^*)$ to $\Gamma(P)$ and then mapping to $\Gamma(P)^*$ (the last isomorphism is due to it being
    fg free).
    This is just $\phi$, so in the case that $P$ is free we have that $\phi$ is an isomorphism.

    For general $E$, by Serre-Swan there is an $F$ such that $E\oplus F\cong P$ trivial so we have $P^*\cong(E\oplus F)^*\cong E^*\oplus F^*$ where the last isomorphism is due to the
    lifting of the natural isomorphism of vector spaces.
    Thus we have an isomorphism (the first isomorphism is due to the above lemma)
    $$ \Gamma(E)^*\oplus\Gamma(F)^* \cong \Gamma(P)^* \cong \Gamma(P^*) \cong \Gamma(E^*\oplus F^*) \cong \Gamma(E^*)\oplus\Gamma(F^*) $$
    Then we have a commutative diagram

    \commsq{\Gamma(E^*)&\Gamma(E^*)\oplus(F^*)\cr \Gamma(E)^*&\Gamma(E)^*\oplus\Gamma(F)^*}
        {}{}{\phi}{\cong}

    where the horizontal arrows are inclusions.
    This shows that $\phi$ is injective, and reversing the horizontal arrows to give projections shows $\phi$ is surjective.
    Note that this is the same as saying the map between free modules is the direct sum of the maps between projective modules.
    \qed

\eproof

We can also find $\phi$'s inverse explicitly.
Let $\psi\colon\Gamma(E)^*\to\Gamma(E^*)$ by mapping $f\colon\Gamma(E)\to C^\infty(M)$ to $\psi(f)\in\Gamma(E^*)$ given by $\psi(f)_p\in E^*_p$
which maps $v\in E_p$ to $f(\bar v)(p)$ where $\bar v$ is a global section extending $v$, then $\psi=\phi^{-1}$.
The above proof though generalizes nicely to prove a similar result for tensor products.

\bprop

    For vector bundles $E,F$,
    $$ \Gamma(E\otimes F)\cong\Gamma(E)\otimes\Gamma(F), $$
    where the latter tensor product is taken over $C^\infty(M)$.

\eprop

\bproof

    There is a natural map $\Gamma(E)\otimes\Gamma(F)\to\Gamma(E\otimes F)$ given by mapping $s\otimes t$ to $s\otimes t\in\Gamma(E\otimes F)$ defined by
    $(s\otimes t)|_p=s|_p\otimes t|_p$.
    This is easily shown to be an isomorphism when $E,F$ are trivial this is an isomorphism.
    Then in the general case we write $E,F$ as direct summands of trivial bundles and use a similar argument as above to show that this map is an isomorphism
    in general.
    \qed

\eproof

\bcoro

    There is a natural isomorphism of vector bundles $(E\otimes F)^*\cong E^*\otimes F^*$.

\ecoro

\bproof

    There is an isomorphism
    $$ \Gamma((E\otimes F)^*) \cong \Gamma(E\otimes F)^* \cong \Gamma(E^*)\otimes\Gamma(F^*) \cong \Gamma(E^*\otimes F^*) , $$
    and by Serre-Swan since $\Gamma$ is an equivalence we have $(E\otimes F)^*\cong E^*\otimes F^*$.
    \qed

\eproof

\bcoro

    For vector bundles $E,F$,
    $$ \Gamma(\hom(E,F)) \cong \hom_{C^\infty(M)}(\Gamma(E),\Gamma(F)) \cong \hom_{\VBun_M}(E,F) . $$
    Furthermore, for any vector bundle $E$, there is an adjunction $\bul\otimes E\dashv\hom(E,\bul)$.

\ecoro

\bproof

    There is a natural isomorphism $\hom(E,F)\cong E^*\otimes F$, thus
    $$ \Gamma(\hom(E,F)) \cong \Gamma(E^*\otimes F) \cong \Gamma(E)^*\otimes\Gamma(F) \cong \hom_{C^\infty(M)}(\Gamma(E),\Gamma(F)) , $$
    where the final isomorphism is due to the space of sections being fg projective.
    And since $\Gamma$ is fully faithful we have the final isomorphism to vector bundle morphisms.

    The tensor-hom adjunction for vector spaces gives an isomorphism
    $$ \hom(\bul\otimes E,\bul) \cong \hom(\bul,\hom(E,\bul)) . $$
    Then applying $\Gamma$ and appealing to what we just proved, this gives an isomorphism
    $$ \hom_{\VBun_M}(\bul\otimes E,\bul) \cong \hom_{\VBun_M}(\bul,\hom(E,\bul)) , $$
    demonstrating the desired adjunction.
    \qed

\eproof

Recall that we categorized tensor fields as multilinear maps out of $\X^*(M)^k\times\X(M)^\ell$.
We can generalize this easily.

\bthrm

    There is a natural isomorphism between sections of $E\otimes F$ and $C^\infty(M)$-bilinear maps $\Gamma(E)^*\times\Gamma(F)^*\to C^\infty(M)$.
    Inductively, there is a natural isomorphism
    $$ \Gamma(E_1\otimes\cdots\otimes E_n) \cong {\rm MultLin}(\Gamma(E_1)^*,\dots,\Gamma(E_n)^*;\ C^\infty(M)) . $$
    And in particular,
    $$ \Gamma(E^{\otimes k}\otimes (E^*)^{\otimes\ell}) \cong {\rm MultLin}(\Gamma(E^*)^k\times\Gamma(E)^\ell;\ C^\infty(M)) . $$

\ethrm

\bproof

    Since fg projective modules are reflective, there is a natural isomorphism
    $$ \Gamma(E\otimes F) \cong \Gamma(E\otimes F)^{**} \cong \Gamma(E)^*\otimes\Gamma(F)^* = \hom_{C^\infty(M)}(\Gamma(E)^*\otimes\Gamma(F)^*,C^\infty(M)) , $$
    and by the universal property of tensor products this is isomorphic to the space of bilinear maps $\Gamma(E)^*\times\Gamma(F)^*\to C^\infty(M)$.
    The inductive case is then clear, and for the final isomorphism we just utilize the reflexivity of vector bundles: $E^{**}\cong E$.
    \qed

\eproof

